\subsection{SOLVABLE LIE GROUPS}

PROPOSITION 19. Let G be a finite-dimensional Lie group. For L(G) to be solvable, it is necessary and sufficient that G possess a solvable open subgroup.

The proof is analogous to that of Proposition 12 of no. 5.

PROPOSITION 20. Let G be a simply connected solvable Lie group of finite-dimension n over R or C and g = L(G). Let \((\mathfrak{g}_{n}, \mathfrak{g}_{n-1}, \ldots, \mathfrak{g}_{0})\) be a sequence of subalgebras of g of dimensions n, n-1, \ldots, 0, such that \(g_{i-1}\) is an ideal of \(g_{i}\) for i = n, n-1, \ldots, 1.† Let \(G_{i}\) be the integral subgroup of G corresponding to \(g_{i}\) . Let \(x_{i}\) be a vector of \(g_{i}\) not belonging to \(g_{i-1}\) . Let \(\phi_{i}\) be the mapping

\[
\left(\lambda_ {1}, \lambda_ {2}, \dots , \lambda_ {i}\right) \mapsto \left(\exp \lambda_ {1} x _ {1}\right) \left(\exp \lambda_ {2} x _ {2}\right) \dots \left(\exp \lambda_ {i} x _ {i}\right)
\]

of \(\mathbf{K}^i\) into G. Then \(\phi_n\) is an isomorphism of analytic manifolds and \(\phi_i(\mathbf{K}^i) = \mathbf{G}_i\) for all \(i\) .

For \(n = 0\) the proposition is obvious. We argue by induction on \(n\) . Let H be the integral subgroup of G such that \(\mathrm{L(H)} = \mathrm{K}x_n\) . By § 6, no. 6, Corollary 1 to Proposition 14, H and \(G_{n-1}\) are simply connected Lie subgroups of G and, as a Lie group, G is the semi-direct product of H by \(G_{n-1}\) . Hence \(\lambda \mapsto \exp (\lambda x_n)\) is an isomorphism of K onto H and by the induction hypothesis the mapping

\[
\left(\lambda_ {1}, \lambda_ {2}, \dots , \lambda_ {n - 1}\right) \mapsto \left(\exp \lambda_ {1} x _ {1}\right) \left(\exp \lambda_ {2} x _ {2}\right) \dots \left(\exp \lambda_ {n - 1} x _ {n - 1}\right)
\]

is an isomorphism of the analytic manifold \(\mathbf{K}^{n - 1}\) onto the analytic manifold \(\mathbf{G}_{n - 1}\) which maps \(\mathbf{K}^i\times \{0\}\) to \(\mathbf{G}_i\) for \(i = 1,2,\ldots ,n - 1\) . Hence the proposition.

PROPOSITION 21. Let G be a simply connected finite-dimensional solvable Lie group over R or C and M an integral subgroup of G. Then M is a Lie subgroup of G and is simply connected.

We continue to use the notation \(n\) , \(\mathfrak{g}\) , \(\mathfrak{g}_{i}\) , \(x_{i}\) , \(\phi\) of Proposition 20 but impose on the \(x_{i}\) the following supplementary condition: let \(i_{p} > i_{p-1} > \cdots > i_{1}\) be the integers \(i\) such that \(\mathrm{L}(M) \cap \mathfrak{g}_{i} \neq \mathrm{L}(M) \cap \mathfrak{g}_{i-1}\) ; then we take

\[
x _ {i _ {k}} \in \mathrm{L} (\mathrm{M}) \cap \mathfrak {g} _ {i _ {k}}
\]

for \(k = 1, 2, \ldots, p\) . By induction on \(n\) , it is easily seen that \((x_{i_p}, x_{i_{p-1}}, \ldots, x_{i_1})\) is a basis of \(\mathrm{L}(\mathbf{M})\) . Let \(\mathbf{N}\) be a simply connected Lie group such that there exists an isomorphism \(h\) of \(\mathrm{L}(\mathbf{N})\) onto \(\mathrm{L}(\mathbf{M})\) . Let

\[
y _ {p} = h ^ {- 1} (x _ {i _ {p}}), \dots , y _ {1} = h ^ {- 1} (x _ {i _ {1}}).
\]

By Proposition 20, the mapping

\[
\left(\lambda_ {1}, \lambda_ {2}, \dots , \lambda_ {p}\right) \mapsto \left(\exp \lambda_ {1} y _ {1}\right) \left(\exp \lambda_ {2} y _ {2}\right) \dots \left(\exp \lambda_ {p} y _ {p}\right)
\]

is an isomorphism of the manifold \(K^{p}\) onto the manifold N. There exists a Lie group morphism \(\tau\) of N into G such that \(h = L(\tau)\) and \(\tau(N) = M\) ( \(\S 6\) , no. 2, Corollary 1 to Proposition 1). Hence M is the set of elements of G of the form

\[
\begin{array}{r l} \tau ((\exp \lambda_ {1} y _ {1}) \dots (\exp \lambda_ {p} y _ {p})) & = \exp (\lambda_ {1} L (\tau) y _ {1}) \dots \exp (\lambda_ {p} L (\tau) y _ {p}) \\ & = \exp (\lambda_ {1} x _ {i _ {1}}) \dots \exp (\lambda_ {p} x _ {i _ {p}}). \end{array}
\]

Thus \(\mathbf{M} = \phi (\mathbf{T})\) where \(\mathbf{T}\) is a vector subspace of \(\mathbf{K}^n\) .

PROPOSITION 22. Suppose that K = R or C. Let V be a finite-dimensional vector space and G a connected solvable subgroup of GL(V). Suppose that the identity representation of G is simple.

\begin{enumerate}
\def\labelenumi{(\roman{enumi})}
\item
  If \(\mathbf{K} = \mathbf{R}\) , then \(\dim V \leqslant 2\) and \(G\) is commutative.
\item
  If K = C, then dim V = 1.
\item
  Suppose that \(\mathbf{K} = \mathbf{R}\) . Then the closure H of G in \(\mathbf{GL}(V)\) is a solvable connected Lie subgroup of \(\mathbf{GL}(V)\) (no. 1, Corollary 2 to Proposition 1). Hence L(H) is solvable (Proposition 19). The identity representation of L(G) is simple ( \(\S 6\) , no. 5, Corollary 2 to Proposition 13). Hence dim V \(\leqslant 2\) and L(G) is commutative (Chapter I, \(\S 5\) , Corollaries 1 and 4 to Theorem 1). Hence G is commutative.
\item
  Suppose that K = C. Let W be a minimal element among the nonzero real vector subspaces of V which are stable under G. The complex vector subspace of V generated by W is equal to V since the identity representation of G is simple. By (i), G\textbar W is commutative. Hence G is commutative. Therefore every element of G is a homothety (Algebra, Chapter VIII, § 4, Corollary I to Proposition 2), so that dim V = 1.
\end{enumerate}

COROLLARY. Let V be a complex vector space of finite dimension \textgreater0 and G a connected solvable subgroup of GL(V).

\begin{enumerate}
\def\labelenumi{(\roman{enumi})}
\item
  There exists a non-zero element \(v\) of \(V\) such that \(gv \in \mathbf{C}v\) for all \(g \in G\) .
\item
  There exists a basis \(\mathbf{B}\) of \(\mathbf{V}\) such that, for all \(g \in \mathbf{G}\) , the matrix of \(g\) with respect to \(\mathbf{B}\) is lower triangular.
\end{enumerate}

Let \(V_{1}\) be a minimal element among the non-zero vector subspaces of V which are stable under G. By Proposition 22 (ii), \(\dim V_{1} = 1\) . This proves (i). By induction on \(\dim V\) , it then follows that there exists an increasing sequence \((\mathrm{V}_1, \mathrm{V}_2, \ldots, \mathrm{V}_n)\) of vector subspaces of V which are stable under G such that \(\dim \mathrm{V}_{i+1} / \mathrm{V}_i = 1\) for \(i < n\) and \(\mathrm{V}_n = \mathrm{V}\) ; hence (ii).

